% بنك التمارين — الفصل 4
% يُحدَّد الفصل من اسم هذا الملف.

\begin{exo}[تكامل شكل تفاضلي على طول مسار]{integrale-forme-differentielle}
\keywords{شكل تفاضلي، تفاضل تام، معيار شوارز، تكامل خطي}

لنعتبر الشكلين التفاضليين
\[
\omega_1 = y\,\mathrm{d}x + x\,\mathrm{d}y,
\qquad
\omega_2 = y\,\mathrm{d}x,
\]
وثلاثة مسارات تصل المبدأ $O(0,0)$ بالنقطة $M(1,1)$: يتكون المسار
$\gamma_1$ من القطعة الأفقية $O\to(1,0)$ تليها القطعة العمودية
$(1,0)\to M$؛ ويتكون المسار $\gamma_2$ من القطعة العمودية $O\to(0,1)$
تليها القطعة الأفقية $(0,1)\to M$؛ أما المسار $\gamma_3$ فهو القطعة المستقيمة
التي تصل $O$ مباشرةً بـ $M$.

\begin{enumerate}
\item بيّن أن $\omega_1$ تام بإيجاد دالة $f$ تحقق $\omega_1=\mathrm{d}f$.
\item احسب $\int_{\gamma_1}\omega_1$ و$\int_{\gamma_2}\omega_1$ بتمثيل كل قطعة وسيطيًا. استنتج النتيجة من دون حساب.
\item طبّق معيار شوارز على $\omega_2$. ماذا نستنتج؟
\item احسب $\int_{\gamma_1}\omega_2$ و$\int_{\gamma_2}\omega_2$ و$\int_{\gamma_3}\omega_2$، ثم علّق.
\end{enumerate}

\begin{indication}
على قطعة أفقية $\mathrm{d}y=0$، وعلى قطعة عمودية $\mathrm{d}x=0$؛ لذلك
يُختزل كل تكامل إلى تكامل بسيط.
\end{indication}

\begin{solution}
\textbf{السؤال 1.} تصلح $f(x,y)=xy$ لأن $\partial f/\partial x=y$ و$\partial f/\partial y=x$.

\textbf{السؤال 2.} لمنحنى ممثّل وسيطيًا بالعلاقة
$t\mapsto\bigl(x(t),y(t)\bigr)$ لدينا
\[
\int_\gamma\omega_1
= \int \left[y(t)x'(t)+x(t)y'(t)\right]\,\mathrm{d}t.
\]

المسار $\gamma_1$ اتحاد قطعتين. على الأولى $x(t)=t$ و$y(t)=0$، حيث
$t\in[0,1]$؛ ومن ثم $\mathrm{d}x=\mathrm{d}t$ و$\mathrm{d}y=0$.
وعلى الثانية $x(t)=1$ و$y(t)=t$، مع $t\in[0,1]$ أيضًا؛ ومن ثم
$\mathrm{d}x=0$ و$\mathrm{d}y=\mathrm{d}t$. لذلك
\[
\begin{aligned}
I_1
= \int_{\gamma_1}\omega_1
&= \int_0^1\left(0\times 1+t\times 0\right)\,\mathrm{d}t
 + \int_0^1\left(t\times 0+1\times 1\right)\,\mathrm{d}t \\
&= 0+\int_0^1\mathrm{d}t = 1.
\end{aligned}
\]

في $\gamma_2$ تُمثّل القطعة الأولى بـ $x(t)=0$ و$y(t)=t$، ثم الثانية بـ
$x(t)=t$ و$y(t)=1$، حيث $t\in[0,1]$ في الحالتين. إذن
\[
\begin{aligned}
I_2
= \int_{\gamma_2}\omega_1
&= \int_0^1\left(t\times 0+0\times 1\right)\,\mathrm{d}t
 + \int_0^1\left(1\times 1+t\times 0\right)\,\mathrm{d}t \\
&= 0+\int_0^1\mathrm{d}t = 1.
\end{aligned}
\]
القيمتان متساويتان. وكان يمكن توقّع ذلك من دون حساب: بما أن
$\omega_1=\mathrm{d}f$ مع $f(x,y)=xy$، فإن تكامله لا يعتمد إلا على الطرفين،
\[
\int_\gamma\omega_1=f(M)-f(O)=f(1,1)-f(0,0)=1.
\]

\textbf{السؤال 3.} نكتب $\omega_2=A(x,y)\,\mathrm{d}x+B(x,y)\,\mathrm{d}y$.
هنا $A(x,y)=y$ و$B(x,y)=0$. لو كان الشكل تامًا لأعطى معيار شوارز
\[
\frac{\partial A}{\partial y}=\frac{\partial B}{\partial x}.
\]
لكن $\partial A/\partial y=1$ في حين أن $\partial B/\partial x=0$.
المشتقتان المتصالبتان مختلفتان، ولذلك ليس $\omega_2$ تامًا. وقد يعتمد تكامله
على المسار المتّبع بين $O$ و$M$، كما يؤكد الحساب الآتي.

\textbf{السؤال 4.} بما أن $\omega_2=y\,\mathrm{d}x$، فلا يسهم في التكامل إلا
تغير $x$ عند قيمة غير صفرية لـ $y$. باستعمال التمثيلات السابقة نجد على $\gamma_1$
\[
\begin{aligned}
J_1
=\int_{\gamma_1}\omega_2
&=\int_0^1 0\times 1\,\mathrm{d}t
  +\int_0^1 t\times 0\,\mathrm{d}t \\
&=0.
\end{aligned}
\]
على القطعة الأولى من $\gamma_2$ لدينا $x(t)=0$، ومن ثم $\mathrm{d}x=0$؛
وعلى الثانية $x(t)=t$ و$y(t)=1$ و$\mathrm{d}x=\mathrm{d}t$. لذلك
\[
\begin{aligned}
J_2
=\int_{\gamma_2}\omega_2
&=\int_0^1 t\times 0\,\mathrm{d}t
  +\int_0^1 1\times 1\,\mathrm{d}t \\
&=1.
\end{aligned}
\]
وأخيرًا يُمثَّل القطر $\gamma_3$ صراحةً بـ
\[
x(t)=t,\qquad y(t)=t,\qquad t\in[0,1],
\]
ومن ثم $\mathrm{d}x=\mathrm{d}t$ و$\mathrm{d}y=\mathrm{d}t$. فنحصل على
\[
J_3
=\int_{\gamma_3}\omega_2
=\int_0^1 y(t)x'(t)\,\mathrm{d}t
=\int_0^1 t\,\mathrm{d}t
=\frac12.
\]
للمسارات الثلاثة الطرفان نفسيهما، لكنها تعطي ثلاث قيم مختلفة:
$J_1=0$ و$J_2=1$ و$J_3=1/2$. وهذا بالضبط اعتماد الشكل التفاضلي غير التام على المسار.
\end{solution}
\end{exo}

\begin{exo}[التغير نفسه في الطاقة الداخلية بثلاث طرائق انتقال]{meme-delta-u-trois-transferts}
\keywords{المبدأ الأول، الطاقة الداخلية، الحرارة، الشغل الميكانيكي، الشغل الكهربائي، المسعرية، دالة حالة}

يشكّل سائل ومسعره ومقاومة صغيرة مغمورة نظامًا مغلقًا ساكنًا عيانيًا. الوعاء
صلب، وتُفترض السعة الحرارية الكلية للنظام ثابتة:
\[
C=2{,}00\ \mathrm{kJ\,K^{-1}}.
\]
نريد نقل النظام من حالة التوازن $A$ عند $T_A=20{,}0\,^\circ\mathrm{C}$
إلى حالة التوازن $B$ عند $T_B=25{,}0\,^\circ\mathrm{C}$. ونقبل أن
\[
U(B)-U(A)=C(T_B-T_A).
\]
تغيرا طاقتي الحركة والوضع العيانيتين معدومان، وتُهمل جميع الضياعات نحو الخارج.

تُنفَّذ البروتوكولات الثلاثة الآتية كلٌّ على حدة انطلاقًا من الحالة الابتدائية نفسها $A$:
\begin{itemize}
\item[البروتوكول 1] يوضع النظام في تماس حراري مع منبع ساخن من دون تزويده بأي شغل.
\item[البروتوكول 2] الجدران عازلة للحرارة. تحرّك ريشة السائل بفضل سقوط كتلة $m$ مسافة $h=5{,}00$ m. تستقر الريشة والسائل في النهاية، وتنتقل إلى النظام كل الطاقة الميكانيكية التي تفقدها الكتلة. نأخذ $g=9{,}81\ \mathrm{m\,s^{-2}}$.
\item[البروتوكول 3] يستقبل النظام حرارة $Q_3=4{,}00$ kJ، وتُزوَّد المقاومة ببقية الطاقة كهربائيًا. القدرة الكهربائية المستقبلة ثابتة وتساوي $\mathcal P=300$ W.
\end{itemize}

\begin{enumerate}
\item احسب تغير الطاقة الداخلية $\Delta U=U(B)-U(A)$ المشترك بين البروتوكولات الثلاثة.
\item لكل من البروتوكولين الأولين، حدّد الحرارة $Q$ والشغل $W$ اللذين يستقبلهما النظام. وفي البروتوكول الثاني احسب الكتلة $m$ اللازمة.
\item في البروتوكول الثالث، احسب الشغل الكهربائي المستقبَل ومدة تغذية المقاومة.
\item للتحولات الثلاثة حالتا البداية والنهاية نفسيهما. فسّر لماذا قد تختلف مع ذلك قيمتا $Q$ و$W$. هل يحتوي النظام «حرارة» أو «شغلًا» في الحالة النهائية $B$؟
\end{enumerate}

\begin{indication}
طبّق $\Delta U=Q+W$ وفق اصطلاح الإشارة المصرفي. الشغل الذي تستقبله الريشة
يساوي نقصان طاقة وضع الكتلة بمقدار $mgh$. وتُحسب الطاقة الكهربائية العابرة
للحدود عبر الأسلاك شغلًا كهربائيًا.
\end{indication}

\begin{solution}
\textbf{السؤال 1.} فرق درجة الحرارة هو $T_B-T_A=5{,}0$ K. لذلك
\[
\Delta U=C(T_B-T_A)
=2{,}00\times5{,}0
=10{,}0\ \mathrm{kJ}.
\]
لا تعتمد هذه القيمة إلا على حالتي التوازن $A$ و$B$.

\textbf{السؤال 2.} في البروتوكول الأول لا يستقبل النظام أي شغل: $W_1=0$.
ومن ثم يعطي المبدأ الأول
\[
Q_1=\Delta U=10{,}0\ \mathrm{kJ}.
\]
الحرارة موجبة لأن الطاقة تدخل النظام.

في البروتوكول الثاني الجدران عازلة للحرارة، ومن ثم $Q_2=0$. يأتي تغير الطاقة
الداخلية كله من الشغل الميكانيكي للريشة:
\[
W_2=\Delta U=10{,}0\ \mathrm{kJ}.
\]
وبما أن $W_2=mgh$، نجد
\[
m=\frac{W_2}{gh}
=\frac{10{,}0\times10^3}{9{,}81\times5{,}00}
\simeq 204\ \mathrm{kg}.
\]
تذكّرنا هذه القيمة الكبيرة بأن ارتفاعًا متواضعًا لدرجة الحرارة يقابل طاقة
ميكانيكية معتبرة.

\textbf{السؤال 3.} يفرض المبدأ الأول
\[
W_3=\Delta U-Q_3=10{,}0-4{,}00=6{,}00\ \mathrm{kJ}.
\]
هذا الشغل موجب: التغذية الكهربائية تمنح النظام طاقة. ومع $W_3=\mathcal P\,\Delta t$،
\[
\Delta t=\frac{W_3}{\mathcal P}
=\frac{6{,}00\times10^3}{300}
=20{,}0\ \mathrm{s}.
\]

\textbf{السؤال 4.} تعطي المسارات الثلاثة على الترتيب
\[
(Q_1,W_1)=(10{,}0,0)\ \mathrm{kJ},
\qquad
(Q_2,W_2)=(0,10{,}0)\ \mathrm{kJ},
\]
و
\[
(Q_3,W_3)=(4{,}00,6{,}00)\ \mathrm{kJ}.
\]
في جميع الحالات يساوي المجموع $Q+W$ القيمة $10{,}0$ kJ وينتج التغير نفسه
$\Delta U$. الطاقة الداخلية دالة حالة، في حين تصف الحرارة والشغل انتقالات
الطاقة خلال تحول. في الحالة النهائية $B$ يملك النظام طاقة داخلية، لكنه لا
يحتوي «حرارة» ولا «شغلًا».
\end{solution}
\end{exo}

\begin{exo}[انضغاط تحت ضغط خارجي ثابت]{compression-pression-exterieure-constante}
\seoready{true}
\keywords{شغل قوى الضغط، ضغط خارجي ثابت، انضغاط، تمدد، تمدد في الفراغ}

يُضغط غاز من الحجم $V_A$ إلى الحجم $V_B<V_A$ تحت ضغط خارجي ثابت $P_0$.

\begin{enumerate}
\item احسب الشغل الذي يستقبله الغاز.
\item يعود الغاز من $V_B$ إلى $V_A$ في مواجهة الضغط الخارجي نفسه $P_0$. احسب الشغل المستقبَل وقارنه بشغل الانضغاط.
\item أعد دراسة التمدد من $V_B$ إلى $V_A$ عندما يتمدد الغاز في الفراغ. فسّر الإشارات الثلاث الناتجة.
\end{enumerate}

\begin{indication}
استخدم $W=-\int P_{\mathrm{ext}}\,\mathrm{d}V$. يجب تكامل الضغط
\emph{الخارجي}، حتى عندما لا يكون التحول شبه ساكن.
\end{indication}

\begin{solution}
\textbf{السؤال 1.} الضغط الخارجي ثابت، لذلك
\[
W_{A\to B}=-P_0(V_B-V_A)=P_0(V_A-V_B)>0.
\]
يستقبل الغاز شغلًا أثناء الانضغاط.

\textbf{السؤال 2.} في العودة،
\[
W_{B\to A}=-P_0(V_A-V_B)<0.
\]
الشغلان متعاكسان لأن التحولين يُنفّذان في مواجهة الضغط الخارجي الثابت نفسه.

\textbf{السؤال 3.} في الفراغ $P_{\mathrm{ext}}=0$، ومن ثم
\[
W_{B\to A}^{\mathrm{vide}}=0.
\]
إذن لا يعني التمدد تلقائيًا شغلًا سالبًا؛ بل يجب أن تُدفَع قوة خارجية فعلًا.
\end{solution}
\end{exo}
